(1) 带正电 \(q\)、质量 \(m\) 的粒子处于匀强磁场 \(\mathbf B=B\hat z\) 和静电四极势 \(\Phi=V_0[z^2-(x^2+y^2)/2]/(2d^2)\)，其中 \(B,V_0,d>0\)。推导三维牛顿—洛伦兹运动方程。
(2) 求轴向简谐频率 \(\omega_z\)，以复坐标 \(u=x+iy\) 求横向两支正的回转频率 \(\omega_\pm\)。
(3) 推导横向频率均为实数的阱稳定条件，证明 \(\omega_++\omega_-=\omega_c\) 及 \(2\omega_+\omega_-=\omega_z^2\)。
(4) 在强磁场极限 \(\omega_z\ll\omega_c\) 下展开两频率，并解释磁控漂移为何明显慢于修正回旋运动。
